% GENERATED FILE - do not hand-edit.
% Source of truth: data/calculus-reviewed/t-10.json in the zenith-website repo.
% Regenerate with: python3 scripts/build_calculus_tex.py
\documentclass[11pt]{article}
\usepackage[utf8]{inputenc}
\usepackage[T1]{fontenc}
\usepackage{amsmath,amssymb}
\usepackage{graphicx}
\usepackage[margin=1in]{geometry}
\graphicspath{{../..//}}
\setlength{\parindent}{0pt}
\setlength{\parskip}{0.6em}
\begin{document}
\section*{Chapter 10: Sequences and Series --- Test T10}

\bigskip\noindent\textbf{T1.}\ Which sequence converges?
\begin{list}{}{\setlength{\leftmargin}{3em}\setlength{\labelwidth}{2.5em}\setlength{\itemsep}{0pt}\setlength{\topsep}{2pt}}
\item[(A)] $a_n = n+\dfrac{5}{n}$
\item[(B)] $a_n = 2+\dfrac{(-1)^n}{n}$
\item[(C)] $a_n = \cos\dfrac{n\pi}{2}$
\item[(D)] $a_n = \dfrac{n!}{4^n}$
\end{list}

\bigskip\noindent\textbf{T2.}\ If $s_n = 2+\dfrac{(-1)^n}{n^2}$, then
\begin{list}{}{\setlength{\leftmargin}{3em}\setlength{\labelwidth}{2.5em}\setlength{\itemsep}{0pt}\setlength{\topsep}{2pt}}
\item[(A)] $s_n$ diverges by oscillation
\item[(B)] $s_n$ converges to zero
\item[(C)] $\displaystyle\lim_{n\to\infty} s_n = 2$
\item[(D)] $s_n$ diverges to infinity
\end{list}

\bigskip\noindent\textbf{T3.}\ The sequence $a_n = \sin\dfrac{n\pi}{6}$
\begin{list}{}{\setlength{\leftmargin}{3em}\setlength{\labelwidth}{2.5em}\setlength{\itemsep}{0pt}\setlength{\topsep}{2pt}}
\item[(A)] is unbounded
\item[(B)] converges to a number less than 1
\item[(C)] is bounded
\item[(D)] diverges to infinity
\end{list}

\bigskip\noindent\textbf{T4.}\ Which of the following sequences diverges?
\begin{list}{}{\setlength{\leftmargin}{3em}\setlength{\labelwidth}{2.5em}\setlength{\itemsep}{0pt}\setlength{\topsep}{2pt}}
\item[(A)] $a_n = \dfrac{(-1)^{n+1}}{2n}$
\item[(B)] $a_n = \dfrac{2^n}{e^n}$
\item[(C)] $a_n = \dfrac{n^3}{e^n}$
\item[(D)] $a_n = \dfrac{n}{\ln(n+1)}$
\end{list}

\bigskip\noindent\textbf{T5.}\ The sequence $\{r^n\}$ converges if and only if
\begin{list}{}{\setlength{\leftmargin}{3em}\setlength{\labelwidth}{2.5em}\setlength{\itemsep}{0pt}\setlength{\topsep}{2pt}}
\item[(A)] $|r|<1$
\item[(B)] $|r|\leq 1$
\item[(C)] $-1<r\leq 1$
\item[(D)] $0<r<1$
\end{list}

\bigskip\noindent\textbf{T6.}\ $\displaystyle\sum u_n$ is a series of constants for which $\displaystyle\lim_{n\to\infty} u_n=0$. Which of the following statements is always true?
\begin{list}{}{\setlength{\leftmargin}{3em}\setlength{\labelwidth}{2.5em}\setlength{\itemsep}{0pt}\setlength{\topsep}{2pt}}
\item[(A)] $\displaystyle\sum u_n$ converges to a finite sum.
\item[(B)] $\displaystyle\sum u_n$ does not diverge to infinity.
\item[(C)] $\displaystyle\sum u_n$ is a positive series.
\item[(D)] None of statements (A), (B), or (C) are always true.
\end{list}

\bigskip\noindent\textbf{T7.}\ Note that $\dfrac{1}{n(n+2)}=\dfrac{1}{2}\left(\dfrac{1}{n}-\dfrac{1}{n+2}\right)$ $(n\geq 1)$. $\displaystyle\sum_{n=1}^\infty \dfrac{1}{n(n+2)}$ equals
\begin{list}{}{\setlength{\leftmargin}{3em}\setlength{\labelwidth}{2.5em}\setlength{\itemsep}{0pt}\setlength{\topsep}{2pt}}
\item[(A)] $1$
\item[(B)] $\dfrac{1}{2}$
\item[(C)] $\dfrac{3}{4}$
\item[(D)] $\infty$
\end{list}

\bigskip\noindent\textbf{T8.}\ The sum of the geometric series $4-\dfrac{4}{3}+\dfrac{4}{9}-\dfrac{4}{27}+\cdots$ is
\begin{list}{}{\setlength{\leftmargin}{3em}\setlength{\labelwidth}{2.5em}\setlength{\itemsep}{0pt}\setlength{\topsep}{2pt}}
\item[(A)] $3$
\item[(B)] $2$
\item[(C)] $\dfrac{9}{4}$
\item[(D)] $6$
\end{list}

\bigskip\noindent\textbf{T9.}\ Which of the following statements about series is true?
\begin{list}{}{\setlength{\leftmargin}{3em}\setlength{\labelwidth}{2.5em}\setlength{\itemsep}{0pt}\setlength{\topsep}{2pt}}
\item[(A)] If $\displaystyle\lim_{n\to\infty} u_n=0$, then $\displaystyle\sum u_n$ converges.
\item[(B)] If $\displaystyle\lim_{n\to\infty} u_n\neq 0$, then $\displaystyle\sum u_n$ diverges.
\item[(C)] If $\displaystyle\sum u_n$ diverges, then $\displaystyle\lim_{n\to\infty} u_n\neq 0$.
\item[(D)] $\displaystyle\sum u_n$ converges if and only if $\displaystyle\lim_{n\to\infty} u_n=0$.
\end{list}

\bigskip\noindent\textbf{T10.}\ Which of the following series diverges?
\begin{list}{}{\setlength{\leftmargin}{3em}\setlength{\labelwidth}{2.5em}\setlength{\itemsep}{0pt}\setlength{\topsep}{2pt}}
\item[(A)] $\displaystyle\sum \dfrac{1}{n^2}$
\item[(B)] $\displaystyle\sum \dfrac{1}{n^2+3n}$
\item[(C)] $\displaystyle\sum \dfrac{n}{n^3+5}$
\item[(D)] $\displaystyle\sum \dfrac{n}{\sqrt{9n^2-1}}$
\end{list}

\bigskip\noindent\textbf{T11.}\ Which of the following series diverges?
\begin{list}{}{\setlength{\leftmargin}{3em}\setlength{\labelwidth}{2.5em}\setlength{\itemsep}{0pt}\setlength{\topsep}{2pt}}
\item[(A)] $4-\dfrac{4}{5}+\dfrac{4}{25}-\cdots$
\item[(B)] $\dfrac{1}{\sqrt{3}}-\dfrac{1}{\sqrt{4}}+\dfrac{1}{\sqrt{5}}-\cdots$
\item[(C)] $\dfrac{1}{9}-\dfrac{1}{16}+\dfrac{1}{25}-\cdots$
\item[(D)] $1-1.2+1.44-1.728+\cdots$
\end{list}

\bigskip\noindent\textbf{T12.}\ Let $S=\displaystyle\sum_{n=1}^\infty \left(\dfrac{3}{4}\right)^n$; then $S$ equals
\begin{list}{}{\setlength{\leftmargin}{3em}\setlength{\labelwidth}{2.5em}\setlength{\itemsep}{0pt}\setlength{\topsep}{2pt}}
\item[(A)] $\dfrac{4}{3}$
\item[(B)] $\dfrac{3}{2}$
\item[(C)] $2$
\item[(D)] $3$
\end{list}

\bigskip\noindent\textbf{T13.}\ Which of the following expansions is impossible?
\begin{list}{}{\setlength{\leftmargin}{3em}\setlength{\labelwidth}{2.5em}\setlength{\itemsep}{0pt}\setlength{\topsep}{2pt}}
\item[(A)] $\sqrt{x-2}$ in powers of $x$
\item[(B)] $\sqrt{x+2}$ in powers of $x$
\item[(C)] $\ln x$ in powers of $(x-1)$
\item[(D)] $\tan x$ in powers of $\left(x-\dfrac{\pi}{4}\right)$
\end{list}

\bigskip\noindent\textbf{T14.}\ The series $\displaystyle\sum_{n=0}^\infty n!(x-4)^n$ converges if and only if
\begin{list}{}{\setlength{\leftmargin}{3em}\setlength{\labelwidth}{2.5em}\setlength{\itemsep}{0pt}\setlength{\topsep}{2pt}}
\item[(A)] $x=0$
\item[(B)] $3<x<5$
\item[(C)] $x=4$
\item[(D)] $3\leq x\leq 5$
\end{list}

\bigskip\noindent\textbf{T15.}\ Let $f(x)=\displaystyle\sum_{n=0}^\infty x^n$. The radius of convergence of $\displaystyle\int_0^x f(t)\,dt$ is
\begin{list}{}{\setlength{\leftmargin}{3em}\setlength{\labelwidth}{2.5em}\setlength{\itemsep}{0pt}\setlength{\topsep}{2pt}}
\item[(A)] $0$
\item[(B)] $1$
\item[(C)] $2$
\item[(D)] $e$
\end{list}

\bigskip\noindent\textbf{T16.}\ The coefficient of $x^4$ in the Maclaurin series for $f(x)=e^{-x/3}$ is
\begin{list}{}{\setlength{\leftmargin}{3em}\setlength{\labelwidth}{2.5em}\setlength{\itemsep}{0pt}\setlength{\topsep}{2pt}}
\item[(A)] $-\dfrac{1}{1944}$
\item[(B)] $\dfrac{1}{1944}$
\item[(C)] $-\dfrac{1}{162}$
\item[(D)] $\dfrac{1}{162}$
\end{list}

\bigskip\noindent\textbf{T17.}\ If an appropriate series is used to evaluate $\displaystyle\int_0^{0.2} x^2 e^{-x^2}\,dx$, then, correct to four decimal places, the definite integral equals
\begin{list}{}{\setlength{\leftmargin}{3em}\setlength{\labelwidth}{2.5em}\setlength{\itemsep}{0pt}\setlength{\topsep}{2pt}}
\item[(A)] $0.0002$
\item[(B)] $0.0013$
\item[(C)] $0.0026$
\item[(D)] $0.0033$
\end{list}

\bigskip\noindent\textbf{T18.}\ If the series $\tan^{-1}1 = 1-\dfrac{1}{3}+\dfrac{1}{5}-\dfrac{1}{7}+\cdots$ is used to approximate $\dfrac{\pi}{4}$ with an error less than $0.002$, then the smallest number of terms guaranteed sufficient by the first-omitted-term alternating-series bound is
\begin{list}{}{\setlength{\leftmargin}{3em}\setlength{\labelwidth}{2.5em}\setlength{\itemsep}{0pt}\setlength{\topsep}{2pt}}
\item[(A)] $100$
\item[(B)] $150$
\item[(C)] $200$
\item[(D)] $250$
\end{list}

\bigskip\noindent\textbf{T19.}\ Let $f(x)=\tan^{-1}x$ and $P_5(x)$ be the Taylor polynomial of degree 5 for $f$ about $x=0$. Given $\tan^{-1}x = \displaystyle\sum_{n=0}^\infty \dfrac{(-1)^n x^{2n+1}}{2n+1}$, it follows that if $-0.5<x<0.5$,
\begin{list}{}{\setlength{\leftmargin}{3em}\setlength{\labelwidth}{2.5em}\setlength{\itemsep}{0pt}\setlength{\topsep}{2pt}}
\item[(A)] $P_5(x)\leq \tan^{-1}x$
\item[(B)] $P_5(x)\geq \tan^{-1}x$
\item[(C)] $P_5(x)>\tan^{-1}x$ if $x<0$ but $P_5(x)<\tan^{-1}x$ if $x>0$
\item[(D)] $P_5(x)<\tan^{-1}x$ if $x<0$ but $P_5(x)>\tan^{-1}x$ if $x>0$
\end{list}

\bigskip\noindent\textbf{T20.}\ Let $f(x)=\tan^{-1}x$ and $P_{11}(x)$ be the Taylor polynomial of degree 11 for $f$ about $x=0$. Given $\tan^{-1}x = \displaystyle\sum_{n=0}^\infty \dfrac{(-1)^n x^{2n+1}}{2n+1}$, it follows that if $-0.5<x<0.5$,
\begin{list}{}{\setlength{\leftmargin}{3em}\setlength{\labelwidth}{2.5em}\setlength{\itemsep}{0pt}\setlength{\topsep}{2pt}}
\item[(A)] $P_{11}(x)\leq \tan^{-1}x$
\item[(B)] $P_{11}(x)\geq \tan^{-1}x$
\item[(C)] $P_{11}(x)>\tan^{-1}x$ if $x<0$ but $P_{11}(x)<\tan^{-1}x$ if $x>0$
\item[(D)] $P_{11}(x)<\tan^{-1}x$ if $x<0$ but $P_{11}(x)>\tan^{-1}x$ if $x>0$
\end{list}

\bigskip\noindent\textbf{T21.}\ Which of the following series converges?
\begin{list}{}{\setlength{\leftmargin}{3em}\setlength{\labelwidth}{2.5em}\setlength{\itemsep}{0pt}\setlength{\topsep}{2pt}}
\item[(A)] $\displaystyle\sum \dfrac{1}{\sqrt[3]{n}}$
\item[(B)] $\displaystyle\sum \dfrac{1}{\sqrt{n}}$
\item[(C)] $\displaystyle\sum \dfrac{1}{8n-3}$
\item[(D)] $\displaystyle\sum \dfrac{3}{n^2-7}$
\end{list}

\bigskip\noindent\textbf{T22.}\ Which of the following series diverges?
\begin{list}{}{\setlength{\leftmargin}{3em}\setlength{\labelwidth}{2.5em}\setlength{\itemsep}{0pt}\setlength{\topsep}{2pt}}
\item[(A)] $\displaystyle\sum_{n=1}^\infty \dfrac{n+2}{n!}$
\item[(B)] $\displaystyle\sum_{n=2}^\infty \dfrac{1}{n\ln n}$
\item[(C)] $\displaystyle\sum_{n=1}^\infty \dfrac{\ln n}{3^n}$
\item[(D)] $\displaystyle\sum_{n=1}^\infty \dfrac{n}{3^n}$
\end{list}

\bigskip\noindent\textbf{T23.}\ For which of the following series does the Ratio Test fail?
\begin{list}{}{\setlength{\leftmargin}{3em}\setlength{\labelwidth}{2.5em}\setlength{\itemsep}{0pt}\setlength{\topsep}{2pt}}
\item[(A)] $\displaystyle\sum \dfrac{1}{n!}$
\item[(B)] $\displaystyle\sum \dfrac{n}{3^n}$
\item[(C)] $1+\dfrac{1}{2^{5/2}}+\dfrac{1}{3^{5/2}}+\dfrac{1}{4^{5/2}}+\cdots$
\item[(D)] $\displaystyle\sum \dfrac{n^n}{n!}$
\end{list}

\bigskip\noindent\textbf{T24.}\ Which of the following alternating series diverges?
\begin{list}{}{\setlength{\leftmargin}{3em}\setlength{\labelwidth}{2.5em}\setlength{\itemsep}{0pt}\setlength{\topsep}{2pt}}
\item[(A)] $\displaystyle\sum \dfrac{(-1)^{n+1}(n+1)}{n+2}$
\item[(B)] $\displaystyle\sum \dfrac{(-1)^{n+1}}{\ln(n+2)}$
\item[(C)] $\displaystyle\sum \dfrac{(-1)^{n-1}}{\sqrt[3]{n}}$
\item[(D)] $\displaystyle\sum \dfrac{(-1)^{n-1}n}{n^2+3}$
\end{list}

\bigskip\noindent\textbf{T25.}\ The power series $x+\dfrac{x^2}{4}+\dfrac{x^3}{9}+\cdots+\dfrac{x^n}{n^2}+\cdots$ converges if and only if
\begin{list}{}{\setlength{\leftmargin}{3em}\setlength{\labelwidth}{2.5em}\setlength{\itemsep}{0pt}\setlength{\topsep}{2pt}}
\item[(A)] $-1<x<1$
\item[(B)] $-1\leq x\leq 1$
\item[(C)] $-1\leq x<1$
\item[(D)] $-1<x\leq 1$
\end{list}

\bigskip\noindent\textbf{T26.}\ The power series $(x+2)-\dfrac{(x+2)^2}{2!}+\dfrac{(x+2)^3}{3!}-\dfrac{(x+2)^4}{4!}+\cdots$ diverges
\begin{list}{}{\setlength{\leftmargin}{3em}\setlength{\labelwidth}{2.5em}\setlength{\itemsep}{0pt}\setlength{\topsep}{2pt}}
\item[(A)] for no real $x$
\item[(B)] if $-3<x\leq -1$
\item[(C)] if $x<-3$ or $x>-1$
\item[(D)] if $-3\leq x<-1$
\end{list}

\bigskip\noindent\textbf{T27.}\ The series obtained by differentiating term by term the series
\[
(x-3)+\dfrac{(x-3)^2}{4}+\dfrac{(x-3)^3}{9}+\dfrac{(x-3)^4}{16}+\cdots
\]
converges for
\begin{list}{}{\setlength{\leftmargin}{3em}\setlength{\labelwidth}{2.5em}\setlength{\itemsep}{0pt}\setlength{\topsep}{2pt}}
\item[(A)] $2\leq x\leq 4$
\item[(B)] $2\leq x<4$
\item[(C)] $2<x\leq 4$
\item[(D)] $1<x<5$
\end{list}

\bigskip\noindent\textbf{T28.}\ The Taylor polynomial of degree 3 at $x=0$ for $f(x)=\sqrt{1-x}$ is
\begin{list}{}{\setlength{\leftmargin}{3em}\setlength{\labelwidth}{2.5em}\setlength{\itemsep}{0pt}\setlength{\topsep}{2pt}}
\item[(A)] $1-\dfrac{x}{2}-\dfrac{x^2}{8}-\dfrac{x^3}{16}$
\item[(B)] $1-\dfrac{x}{2}+\dfrac{x^2}{8}-\dfrac{x^3}{16}$
\item[(C)] $1-\dfrac{x}{2}-\dfrac{x^2}{4}-\dfrac{x^3}{8}$
\item[(D)] $1+\dfrac{x}{2}-\dfrac{x^2}{8}+\dfrac{x^3}{16}$
\end{list}

\bigskip\noindent\textbf{T29.}\ The Taylor polynomial of degree 3 at $x=2$ for $e^x$ is
\begin{list}{}{\setlength{\leftmargin}{3em}\setlength{\labelwidth}{2.5em}\setlength{\itemsep}{0pt}\setlength{\topsep}{2pt}}
\item[(A)] $e^2\left[1+(x-2)+\dfrac{(x-2)^2}{2}+\dfrac{(x-2)^3}{3}\right]$
\item[(B)] $e^2\left[1+(x+2)+\dfrac{(x+2)^2}{2!}+\dfrac{(x+2)^3}{3!}\right]$
\item[(C)] $e^2\left[1+(x-2)+\dfrac{(x-2)^2}{2!}+\dfrac{(x-2)^3}{3!}\right]$
\item[(D)] $e^2\left[1-(x-2)+\dfrac{(x-2)^2}{2!}-\dfrac{(x-2)^3}{3!}\right]$
\end{list}

\bigskip\noindent\textbf{T30.}\ The coefficient of $\left(x-\dfrac{\pi}{3}\right)^3$ in the Taylor series about $\dfrac{\pi}{3}$ of $f(x)=\sin x$ is
\begin{list}{}{\setlength{\leftmargin}{3em}\setlength{\labelwidth}{2.5em}\setlength{\itemsep}{0pt}\setlength{\topsep}{2pt}}
\item[(A)] $-\dfrac{1}{12}$
\item[(B)] $\dfrac{1}{12}$
\item[(C)] $\dfrac{1}{6\sqrt{3}}$
\item[(D)] $-\dfrac{1}{3\sqrt{3}}$
\end{list}

\bigskip\noindent\textbf{T31.}\ The coefficient of $x^2$ in the Maclaurin series for $e^{2\sin x}$ is
\begin{list}{}{\setlength{\leftmargin}{3em}\setlength{\labelwidth}{2.5em}\setlength{\itemsep}{0pt}\setlength{\topsep}{2pt}}
\item[(A)] $0$
\item[(B)] $1$
\item[(C)] $2$
\item[(D)] $4$
\end{list}

\bigskip\noindent\textbf{T32.}\ The coefficient of $(x-1)^4$ in the Taylor series for $x^2\ln x$ about $x=1$ is
\begin{list}{}{\setlength{\leftmargin}{3em}\setlength{\labelwidth}{2.5em}\setlength{\itemsep}{0pt}\setlength{\topsep}{2pt}}
\item[(A)] $-\dfrac{1}{20}$
\item[(B)] $-\dfrac{1}{12}$
\item[(C)] $-\dfrac{1}{5!}$
\item[(D)] $-\dfrac{1}{4!}$
\end{list}

\bigskip\noindent\textbf{T33.}\ The radius of convergence of the series $\displaystyle\sum_{n=1}^\infty \dfrac{x^n}{3^n}\cdot\dfrac{n^n}{n!}$ is
\begin{list}{}{\setlength{\leftmargin}{3em}\setlength{\labelwidth}{2.5em}\setlength{\itemsep}{0pt}\setlength{\topsep}{2pt}}
\item[(A)] $3$
\item[(B)] $\dfrac{3}{e}$
\item[(C)] $\dfrac{e}{3}$
\item[(D)] $\infty$
\end{list}

\bigskip\noindent\textbf{T34.}\ The Taylor polynomial of degree 3 at $x=0$ for $(1+x)^p$, where $p$ is a constant, is
\begin{list}{}{\setlength{\leftmargin}{3em}\setlength{\labelwidth}{2.5em}\setlength{\itemsep}{0pt}\setlength{\topsep}{2pt}}
\item[(A)] $1+px+p(p-1)x^2+p(p-1)(p-2)x^3$
\item[(B)] $1+px+\dfrac{p(p-1)}{2}x^2+\dfrac{p(p-1)(p-2)}{3}x^3$
\item[(C)] $1+px+\dfrac{p(p-1)}{2!}x^2+\dfrac{p(p-1)(p-2)}{3!}x^3$
\item[(D)] $px+\dfrac{p(p-1)}{2!}x^2+\dfrac{p(p-1)(p-2)}{3!}x^3$
\end{list}

\bigskip\noindent\textbf{T35.}\ The Taylor series for $\ln(1+3x)$ about $x=0$ is
\begin{list}{}{\setlength{\leftmargin}{3em}\setlength{\labelwidth}{2.5em}\setlength{\itemsep}{0pt}\setlength{\topsep}{2pt}}
\item[(A)] $3x-\dfrac{9}{2}x^2+9x^3-\dfrac{81}{4}x^4+\cdots$
\item[(B)] $3x-9x^2+27x^3-\cdots$
\item[(C)] $3x-\dfrac{3}{2}x^2+9x^3-\dfrac{81}{4}x^4+\cdots$
\item[(D)] $3x-\dfrac{9}{2}x^2+\dfrac{9}{2}x^3-\dfrac{27}{8}x^4+\cdots$
\end{list}

\bigskip\noindent\textbf{T36.}\ The set of all values of $x$ for which $\displaystyle\sum_{n=1}^\infty \dfrac{n\cdot 3^n}{x^n}$ converges is
\begin{list}{}{\setlength{\leftmargin}{3em}\setlength{\labelwidth}{2.5em}\setlength{\itemsep}{0pt}\setlength{\topsep}{2pt}}
\item[(A)] only $x=0$
\item[(B)] $-3<x<3$
\item[(C)] $|x|>3$
\item[(D)] $|x|\geq 3$
\end{list}

\bigskip\noindent\textbf{T37.}\ The third-degree Taylor polynomial $P_3(x)$ for $\cos x$ about $\dfrac{\pi}{4}$ is
\begin{list}{}{\setlength{\leftmargin}{3em}\setlength{\labelwidth}{2.5em}\setlength{\itemsep}{0pt}\setlength{\topsep}{2pt}}
\item[(A)] $\dfrac{1}{\sqrt{2}}\left(\left(x-\dfrac{\pi}{4}\right)-\dfrac{1}{3!}\left(x-\dfrac{\pi}{4}\right)^3\right)$
\item[(B)] $\dfrac{1}{\sqrt{2}}\left(1+\left(x-\dfrac{\pi}{4}\right)-\dfrac{1}{2}\left(x-\dfrac{\pi}{4}\right)^2+\dfrac{1}{3!}\left(x-\dfrac{\pi}{4}\right)^3\right)$
\item[(C)] $\dfrac{1}{\sqrt{2}}\left(1-\left(x-\dfrac{\pi}{4}\right)-\dfrac{1}{2!}\left(x-\dfrac{\pi}{4}\right)^2+\dfrac{1}{3!}\left(x-\dfrac{\pi}{4}\right)^3\right)$
\item[(D)] $1+\left(x-\dfrac{\pi}{4}\right)-\dfrac{1}{2}\left(x-\dfrac{\pi}{4}\right)^2-\dfrac{1}{6}\left(x-\dfrac{\pi}{4}\right)^3$
\end{list}

\bigskip\noindent\textbf{T38.}\ Let $h$ be a function for which all derivatives exist at $x=1$. If $h(1)=h'(1)=h''(1)=h'''(1)=8$, which third-degree polynomial best approximates $h$ there?
\begin{list}{}{\setlength{\leftmargin}{3em}\setlength{\labelwidth}{2.5em}\setlength{\itemsep}{0pt}\setlength{\topsep}{2pt}}
\item[(A)] $8+8x+8x^2+8x^3$
\item[(B)] $8+8(x-1)+4(x-1)^2+\dfrac{4}{3}(x-1)^3$
\item[(C)] $8+8x+4x^2+x^3$
\item[(D)] $8+8(x-1)+3(x-1)^2+(x-1)^3$
\end{list}

\bigskip\noindent\textbf{TB1.}\ Which of the following statements about series is false?
\begin{list}{}{\setlength{\leftmargin}{3em}\setlength{\labelwidth}{2.5em}\setlength{\itemsep}{0pt}\setlength{\topsep}{2pt}}
\item[(A)] $\displaystyle\sum_{k=1}^\infty u_k = \displaystyle\sum_{k=m}^\infty u_k$, where $m$ is any positive integer.
\item[(B)] If $\displaystyle\sum u_n$ converges, so does $\displaystyle\sum cu_n$ if $c\neq 0$.
\item[(C)] If $\displaystyle\sum a_n$ and $\displaystyle\sum b_n$ converge, so does $\displaystyle\sum(ca_n+b_n)$, where $c\neq 0$.
\item[(D)] Rearranging the terms of a positive convergent series will not affect its convergence or its sum.
\end{list}

\bigskip\noindent\textbf{TB2.}\ Which of the following statements is always true?
\begin{list}{}{\setlength{\leftmargin}{3em}\setlength{\labelwidth}{2.5em}\setlength{\itemsep}{0pt}\setlength{\topsep}{2pt}}
\item[(A)] If $\displaystyle\sum u_n$ converges, then so does the series $\displaystyle\sum |u_n|$.
\item[(B)] If a series is truncated after the $n$th term, then the error is less than the first term omitted.
\item[(C)] If the terms of an alternating series decrease, then the series converges.
\item[(D)] None of statements (A), (B), or (C) are always true.
\end{list}

\bigskip\noindent\textbf{TB3.}\ Which of the following series can be used to compute $\ln 0.7$?
\begin{list}{}{\setlength{\leftmargin}{3em}\setlength{\labelwidth}{2.5em}\setlength{\itemsep}{0pt}\setlength{\topsep}{2pt}}
\item[(A)] $\ln(x-1)$ expanded about $x=0$
\item[(B)] $\ln x$ about $x=0$
\item[(C)] $\ln x$ expanded about $x=1$
\item[(D)] $\ln(x-1)$ expanded about $x=1$
\end{list}

\bigskip\noindent\textbf{TB4.}\ Let $f(x)=\displaystyle\sum_{n=0}^\infty a_nx^n$, $g(x)=\displaystyle\sum_{n=0}^\infty b_nx^n$. Suppose both series converge for $|x|<R$. Let $x_0$ be a number such that $|x_0|<R$. Which of statements (A)--(C) is false?
\begin{list}{}{\setlength{\leftmargin}{3em}\setlength{\labelwidth}{2.5em}\setlength{\itemsep}{0pt}\setlength{\topsep}{2pt}}
\item[(A)] $\displaystyle\sum_{n=0}^\infty (a_n+b_n)(x_0)^n$ converges to $f(x_0)+g(x_0)$.
\item[(B)] $f(x)=\displaystyle\sum_{n=0}^\infty a_nx^n$ is continuous at $x=x_0$.
\item[(C)] $\displaystyle\sum_{n=1}^\infty na_nx_0^{n-1}$ converges to $f'(x_0)$.
\item[(D)] Statements (A)--(C) are all true.
\end{list}

\bigskip\noindent\textbf{TB5.}\ When $\sin x$ is approximated by $x-x^3/3!$ for every $|x|<1.5$ radians, which is the smallest listed number greater than the absolute error throughout that interval?
\begin{list}{}{\setlength{\leftmargin}{3em}\setlength{\labelwidth}{2.5em}\setlength{\itemsep}{0pt}\setlength{\topsep}{2pt}}
\item[(A)] $0.05$
\item[(B)] $0.063$
\item[(C)] $0.08$
\item[(D)] $0.1$
\end{list}

\bigskip\noindent\textbf{TB6.}\ The function $f(x)=\displaystyle\sum_{n=0}^\infty a_nx^n$ and $f'(x)=-f(x)$ for all $x$. If $f(0)=1$, then how many terms of the series are needed to find $f(0.3)$ correct to three decimal places?
\begin{list}{}{\setlength{\leftmargin}{3em}\setlength{\labelwidth}{2.5em}\setlength{\itemsep}{0pt}\setlength{\topsep}{2pt}}
\item[(A)] $2$
\item[(B)] $3$
\item[(C)] $4$
\item[(D)] $5$
\end{list}

\bigskip\noindent\textbf{TB7.}\ The sum of the series $\displaystyle\sum_{n=1}^\infty \left(\dfrac{e^2}{2e}\right)^n$
\begin{list}{}{\setlength{\leftmargin}{3em}\setlength{\labelwidth}{2.5em}\setlength{\itemsep}{0pt}\setlength{\topsep}{2pt}}
\item[(A)] $=1$
\item[(B)] $=\dfrac{2e}{e^2-2e}$
\item[(C)] $=\dfrac{e^2}{2e-e^2}$
\item[(D)] diverges
\end{list}

\bigskip\noindent\textbf{TB8.}\ When $\displaystyle\sum_{n=1}^\infty \dfrac{(-1)^{n-1}}{4n-1}$ is approximated by the sum of its first 300 terms, the error is closest to
\begin{list}{}{\setlength{\leftmargin}{3em}\setlength{\labelwidth}{2.5em}\setlength{\itemsep}{0pt}\setlength{\topsep}{2pt}}
\item[(A)] $0.001$
\item[(B)] $0.002$
\item[(C)] $0.005$
\item[(D)] $0.01$
\end{list}

\bigskip\noindent\textbf{TB9.}\ You wish to estimate $e^x$, over the interval $|x|\leq 3$, with an error less than $0.001$. The Lagrange error term suggests that you use a Taylor polynomial at 0 with degree at least
\begin{list}{}{\setlength{\leftmargin}{3em}\setlength{\labelwidth}{2.5em}\setlength{\itemsep}{0pt}\setlength{\topsep}{2pt}}
\item[(A)] $11$
\item[(B)] $12$
\item[(C)] $13$
\item[(D)] $14$
\end{list}

\end{document}
